% Consul Architecture
% Pretty Good: http://man.hubwiz.com/docset/Consul.docset/Contents/Resources/Documents/docs/internals/architecture.html
% [Consul Architecture](https://developer.hashicorp.com/consul/docs/architecture)
% [Consul Vocabulary](https://developer.hashicorp.com/consul/docs/install/glossary)
% [Consensus Protocol](https://developer.hashicorp.com/consul/docs/architecture/consensus)
% **Fokus: [KVStore](https://developer.hashicorp.com/consul/docs/dynamic-app-config/kv)**
% - 

\subsection{Overview}

\begin{frame}
	\frametitle{General}
	\begin{itemize}
		\item Released in 2014 as a Service Discovery Platform
		\item Implemented in \emph{Go}
		\item Now: DNS, Health Checks, Distributed KV Store, \ldots
		      % DNS to return healthy services
	\end{itemize}
\end{frame}

\begin{frame}
	\frametitle{Communication Types}
	% Two Types: Raft for Consensus, Serf for Gossip
	\begin{columns}[b]
		\begin{column}{0.5\textwidth}
			\begin{block}[Raft Protocol]
				\begin{itemize}
					\item<2-> Cluster State Replication
					\item<3-> Example: KV store, service and node IP addresses, configuration
				\end{itemize}
			\end{block}
		\end{column}
		\begin{column}{0.5\textwidth}
			\begin{block}[Surf Protocol]
				\begin{itemize}
					\item<4-> Distribute and perform \emph{node health checks}
					\item<5-> Examples for health checks: HTTP GET Request, gRPC, TCP, UDP
						% https://developer.hashicorp.com/consul/docs/services/usage/checks
				\end{itemize}
			\end{block}
		\end{column}
	\end{columns}
\end{frame}

\begin{frame}
	\frametitle{Architecture}
	\contourlength{.09em}
	\begin{tikzpicture}[overlay]
		\node[roundednode, xshift=-1.9cm,yshift=-1.5cm, text width=3cm](logo) at
		(current page.north east) {\small
		\textcolor{yellow}{---} Raft consensus\\
		\textcolor{red}{---} KV request flow\\
		\textcolor{violet}{---} LAN gossip\\
		\textcolor{green}{---} WAN gossip\\};
	\end{tikzpicture}
	\centering
	\begin{tikzpicture}
		\begin{scope}
			\node [roundednode] (S1) [] {Server 1};
			\node [roundednode, draw=accent, fill=accent!10] (S2) [right=of S1] {Server 2};
			\node [roundednode] (S3) [right=of S2] {Server 3};

			\node [draw=none, visible on=<2>] (L) [above=of S1, yshift=1cm] {Leader};
			\draw [textarrow, visible on=<2>] (L) to[bend left=20] (S2);

			\foreach \i in {1,...,3}{
				\node [roundednode] (C\i) [below=of S\i, yshift=-5mm] {Client \i};
			}

			\node [draw=none] (DS1) [above=of S2] {\contour{white}{Datacenter A}};
			\node [roundednode, node on layer=background,draw=darkgray, dashed,
			fill=gray!10, inner sep=3mm] [fit=(S3)(C1)(DS1)] {};

			\node [roundednode] (APP) [below=of C3, yshift=-5mm] {Application};

			\begin{scope}[visible on=<6->]
				\node [roundednode] (SN) [left=of S1, xshift=-1cm] {Server n};
				\node [draw=none] (PC) [left=of C1, xshift=-1cm] {\phantom{C}};
				\node [draw=none] (DS2) [above=of SN] {Datacenter B};
				\node [roundednode, visible on=<6->, node on
				layer=background,draw=darkgray, dashed, fill=gray!10, inner sep=2mm] [fit=(SN)(PC)(DS2)] {};
			\end{scope}

			\foreach \i in {1,...,3}{
				\foreach \j in {1,...,3}{
					\draw [doublearrow, dotted, draw=darkgray, visible on=<{-5,6-8}>] (S\i) to (C\j);
					\draw [doublearrow, draw=violet, visible on=<{5,8}>] (S\i) to (C\j);
				}
				\draw [doublearrow, dotted, draw=darkgray, visible on=<6>] (SN) to [bend left=20] (S\i);
				\draw [doublearrow, draw=green, visible on=<7->] (SN) to [bend left=20] (S\i);
			}

			\begin{scope}[transform canvas={yshift=-1mm}]
				\draw [doublearrow, visible on=<{-3, 4-8}>, dotted, draw=darkgray] (S1) to (S2);
				\draw [doublearrow, visible on=<{-3, 4-8}>, dotted, draw=darkgray] (S2) to (S3);
				\draw [doublearrow, visible on=<{3, 8}>, draw=yellow] (S1) to (S2);
				\draw [doublearrow, visible on=<{3, 8}>, draw=yellow] (S2) to (S3);
			\end{scope}

			\begin{scope}[transform canvas={xshift=3mm}]
				\draw [arrow, visible on=<{-4, 5-8}>, dotted, draw=darkgray] (C3) to node [anchor=west] {\contour{white}{RPC}} (S3);
				\draw [arrow, visible on=<{-4, 5-8}>, dotted, draw=darkgray] (APP) to node [anchor=west] {HTTP} (C3);
				\draw [arrow, visible on=<{4,8}>, draw=red] (C3) to node [anchor=west] {\contour{white}{RPC}} (S3);
				\draw [arrow, visible on=<{4,8}>, draw=red] (APP) to node [anchor=west] {HTTP} (C3);
			\end{scope}
			\begin{scope}[transform canvas={yshift=1mm}]
				\draw [arrow, visible on=<{-4, 5-8}>, dotted, draw=darkgray] (S3) to (S2);
				\draw [arrow, visible on=<{4,8}>, draw=red] (S3) to (S2);
			\end{scope}
		\end{scope}
	\end{tikzpicture}
	\let\thefootnote\relax\footnote{Based on \autocite[Consul Architecture]{ConsulDoc23}}
	\addtocounter{footnote}{-1}

	% Agent (Server/Client), Datacenters
	% Clients are Stateless and Forward RPC calls to Server Agents
	% Servers participate in Raft Quorum. Thus, they save cluster state (like KV keys)
	% Health checks are done though gossip and not consul -> less traffic with consul replication
	% There is no replication of KV store between datacenters, so if a key is stored in a different datacenter, request is forwarded
\end{frame}
\begin{frame}
	\frametitle{Raft Consensus}
	\framesubtitle{Short Overview}
	\begin{tikzpicture}[remember picture,overlay]
		\node[xshift=-1.5cm,yshift=-1.5cm](logo) at (current page.north east){%
		\includegraphics[width=3cm]{assets/raft_logo.pdf}};
		\node[draw=none,yshift=0.5cm,text width=3cm] at (logo.south) {\baselineskip=1pt\tiny\url{commons.wikimedia.org/w/index.php?curid=105476108}\par};
	\end{tikzpicture}

	\begin{itemize}
		\item<1-> Consensus algorithm for state machine replication\footnote{See \autocite{raft14}}
		\item<2-> Not Byzantine Fault Tolerant (leader is trusted)
		\item<3-> Crash Tolerance: (N/2) + 1
			\begin{itemize}
				\item 3 Nodes: 1 Crash
				\item 5 Nodes: 2 Crashes
				\item 7 Nodes: 3 Crashes
			\end{itemize}
		\item<4-> Three states
	\end{itemize}
	% Logo Credits:
\end{frame}

\begin{frame}
	\frametitle{Consistency}

	\begin{itemize}
		\item<1-> \emph{Writes}: always sent to leader
		\item<2-> \emph{Reads}: three consistency modes:			\begin{description}[labelsep=0.5em]
				\setlength\itemsep{0.5em}
				\item[default]<3-> \emph{leader leasing}: leader assumes its role for a specific time window and responds without quorum
				\\
				(However: risk of 2 concurrent leaders $\Rightarrow$ stale reads)
				\item[consistent]<4-> leader has to verify its role before responding
				%				that they are still the leader
				\item[stale]<5-> any server agent (leader \& follower) can respond
			\end{description}
	\end{itemize}
\end{frame}

\subsection{Usage}

\begin{frame}[fragile]
	\frametitle{CLI agent}
	\framesubtitle{for KV store}

	\begin{columns}[c]
		\begin{column}{0.45\textwidth}
			\begin{minted}{bash}
$ consul agent -dev

# other shell:
$ consul kv put my/key 123
$ consul kv get my/key
123
			\end{minted}
		\end{column}

		\begin{column}{0.55\textwidth}
			\onslide<1->{\begin{block}[KV Store]
					\begin{itemize}
						\item<1-> For configuration, locks, metadata, ...
						\item<2-> Max value size of 512 KB
						\item<3-> Requests to KV store via CLI or HTTP API
					\end{itemize}
				\end{block}}
			% Metadata response: https://developer.hashicorp.com/consul/api-docs/kv#metadata-response
		\end{column}
	\end{columns}
\end{frame}

\begin{frame}[fragile]
	\frametitle{Long Polling}
	\begin{minted}{bash}
$ curl -v http://localhost:8500/v1/kv/my/key
# ...
< X-Consul-Index: 19
# ...

$ curl -v http://localhost:8500/v1/kv/my/key?index=19
# blocks until value is changed or timeout is reached
# (max. 10 minutes)
	\end{minted}
\end{frame}
